\documentclass[11pt,a4paper]{article}
\usepackage[T1]{fontenc}
\usepackage{lmodern}
\usepackage[margin=28mm]{geometry}
\usepackage{amsmath,amssymb,amsthm,mathtools}
\usepackage{microtype}
\usepackage{aliascnt}
\usepackage[colorlinks=true,linkcolor=blue,citecolor=blue,urlcolor=blue]{hyperref}
\usepackage[nameinlink,noabbrev]{cleveref}
\DeclareMathOperator{\Op}{Op}
\DeclareMathOperator{\spec}{spec}
\DeclareMathOperator{\sgn}{sgn}
\DeclareMathOperator{\Tr}{Tr}
\DeclareMathOperator{\ord}{ord}
\newtheorem{theorem}{Theorem}[section]
\newaliascnt{lemma}{theorem}
\newtheorem{lemma}[lemma]{Lemma}
\aliascntresetthe{lemma}
\newaliascnt{proposition}{theorem}
\newtheorem{proposition}[proposition]{Proposition}
\aliascntresetthe{proposition}
\theoremstyle{remark}
\newaliascnt{remark}{theorem}
\newtheorem{remark}[remark]{Remark}
\aliascntresetthe{remark}
\numberwithin{equation}{section}

\title{A sharp Weyl remainder for constant-index transmission eigenvalues}
\author{Kaiwen Jin\thanks{Email: \href{mailto:yui.ui.siki@gmail.com}{\texttt{yui.ui.siki@gmail.com}}.}}
\date{}
\hypersetup{pdftitle={A sharp Weyl remainder for constant-index transmission eigenvalues},pdfauthor={Kaiwen Jin}}

\begin{document}
\maketitle
\begin{abstract}
We prove a Weyl formula with remainder $O(r^{d-1})$ for the interior
transmission problem with a constant refractive index $n>1$ in an arbitrary
bounded smooth domain in $\mathbb R^d$, $d\ge2$. The count includes every
nonzero complex energy of modulus at most $r^2$, with its algebraic
multiplicity for the original maximal-domain transmission operator.
The leading coefficient is
$\omega_d(2\pi)^{-d}|\Omega|(1+n^{d/2})$.
The argument combines the known Dirichlet-to-Neumann approximation with a
sector estimate for a normalized boundary family. At inverse frequency scale $h$, its trace-class part
has norm $O(h^{1-d})$ at scaled energy height $h$, and the sector estimate controls an
actual logarithm of its determinant. Pole cancellation in a strip then
gives the same order for the contour argument, including the endpoint
contribution at every counting radius. No convexity assumption is used.
\end{abstract}

\section{Introduction}

The interior transmission problem asks when two Helmholtz fields with
different wave speeds can have the same Cauchy data on the boundary.
For a constant refractive index $n>1$, it takes the form
\begin{equation}\label{eq:transmission}
 \Delta u+n\lambda u=0,\qquad \Delta v+\lambda v=0
 \quad\hbox{in }\Omega,\qquad
 u=v,\quad \partial_\nu u=\partial_\nu v
 \quad\hbox{on }\partial\Omega.
\end{equation}
The coupling makes the spectral problem nonselfadjoint. Its high-energy
count is nevertheless expected to be the sum of the two scalar
Dirichlet counts, up to an error of boundary order. The main difficulty
is to control the argument of a meromorphic boundary determinant close
to the real spectrum, where the two Dirichlet spectra may meet the
transmission spectrum.

Petkov and Vodev \cite{PV} relate the transmission count, with algebraic
multiplicity, to the scalar Dirichlet counts through an exact boundary
index formula. Their construction reduces the counting remainder to
estimates for a meromorphic Dirichlet-to-Neumann family. Vodev's
high-frequency approximation \cite{Vdn} supplies, in the isotropic case,
a fixed eigenvalue-free strip in the wave-number plane and a counting
remainder $O_\epsilon(r^{d-1+\epsilon})$ for every $\epsilon>0$.
The boundary-order remainder $O(r^{d-1})$ is the conjectured endpoint
\cite[Section~7, Conjecture~1, (19)]{Vsurvey}. We prove that endpoint
for the constant-index problem \eqref{eq:transmission} on every smooth
bounded domain.

We first specify the realization and the count. All spaces are complex.
Let $\Omega\subset\mathbb R^d$ be bounded and connected, with
$C^\infty$ boundary. On $\mathcal H=L^2(\Omega)\oplus L^2(\Omega)$,
let
\begin{equation}\label{eq:A-domain}
 A(u,v)=(-n^{-1}\Delta u,-\Delta v),\qquad
 \begin{aligned}
 \mathcal D(A)=\{(u,v)\in\mathcal H:\;&\Delta u,\Delta v\in L^2(\Omega),\\
 &\gamma_0u=\gamma_0v,\quad\gamma_1u=\gamma_1v\}.
 \end{aligned}
\end{equation}
Here the Laplacians are distributional and $\gamma_0,\gamma_1$ are
the generalized Dirichlet and normal traces on the maximal Laplace
graph domain. Either common choice of normal gives this same domain.
This is the closed transmission realization used throughout the paper.
Its nonzero eigenvalues are isolated and have finite algebraic
multiplicity. For a small positively oriented circle $C_\lambda$
enclosing one such eigenvalue, set
\begin{equation}\label{eq:count-definition}
 m_T(\lambda)=\operatorname{rank}
 \frac{1}{2\pi i}\int_{C_\lambda}(z-A)^{-1}\,dz,
 \qquad
 N(r)=\sum_{0<|\lambda|\le r^2}m_T(\lambda).
\end{equation}
Thus $N(r)$ counts energies, including nonreal energies, with the Riesz
multiplicity of the original operator. It does not count the two square
roots of an energy separately. Finiteness near zero is established in
\cref{lem:regularity}.

\begin{theorem}\label{thm:main}
Let $d\ge2$, let $\Omega\subset\mathbb R^d$ be a bounded connected
domain with $C^\infty$ boundary, and fix a constant $n>1$.
For the operator \eqref{eq:A-domain} and count
\eqref{eq:count-definition},
\begin{equation}\label{eq:main}
 N(r)=\frac{\omega_d}{(2\pi)^d}|\Omega|(1+n^{d/2})r^d
       +O_{\Omega,n}(r^{d-1})\qquad(r\longrightarrow\infty).
\end{equation}
Here $\omega_d$ is the volume of the unit ball in $\mathbb R^d$ and
$|\Omega|$ is the volume of $\Omega$. The estimate holds for every
sufficiently large real radius $r$, including radii through
eigenvalues.
\end{theorem}

The boundary may have several components; no convexity or other
geometric restriction is imposed. The index $n$ is fixed, and the
constant in the remainder may deteriorate as $n\downarrow1$.
The theorem is the constant-index isotropic case of the conjecture
cited above. It makes no assertion for variable refractive indices or
anisotropic principal parts.

The proof retains the Poisson construction and the exact Riesz/index
identity of \cite{PV}, and the approximation of \cite{Vdn}. The
additional estimate concerns the determinant's phase at the critical
height $|\Im\zeta|\asymp h$ in the scaled energy $\zeta=h^2\lambda$,
where $h\asymp r^{-1}$ is the inverse counting scale. After multiplication by one tangential
derivative and a scalar factor, the difference of the two
Dirichlet-to-Neumann maps has principal symbol
\[
 \frac{2\sqrt{1+r_0}}{
 \sqrt{r_0-n\zeta}+\sqrt{r_0-\zeta}},
\]
with square roots of positive real part in the upper half-plane.
Here $r_0$ is the symbol of the positive boundary Laplacian.
This symbol stays in a fixed sector and tends to one at high tangential
frequency. The full approximation in its one-derivative output norm
therefore gives a uniformly coercive boundary family at height $Ch$.
A holomorphic high-frequency normalization leaves a trace-class
perturbation of the identity of size $O(h^{1-d})$. Coercivity along the
straight-line homotopy to the identity bounds an analytic logarithm
of its determinant, including its imaginary part. This is the
phase estimate needed at the endpoint.

The determinant still has Dirichlet poles inside the thin strip.
We cancel their divisors with finite strip factors and use a harmonic
barrier to transfer the horizontal bound into the strip. Jensen's
formula bounds the local transmission multiplicity by $O(h^{1-d})$;
a disk factorization then bounds the argument on each short vertical
side by the same order. The index formula gives the annular count,
and the local estimate removes contour shifts at exact endpoints.
Dyadic summation yields \eqref{eq:main}.

\Cref{sec:prelim} records the realization, scaling and external boundary
estimates. \Cref{sec:det} constructs the normalized determinant and its
logarithm, and \cref{sec:strip} proves the thin-strip estimates.
\Cref{sec:count} completes the count. The radial calculation in
\cref{app:ball}, related to the constant-index models in \cite{PS},
illustrates why common Dirichlet values must retain their original
algebraic multiplicity.


\section{Spectral and boundary preliminaries}\label{sec:prelim}

Throughout the proof, $d$, $\Omega$ and $n$ satisfy the hypotheses of
\cref{thm:main}, and $q=d-1$. Constants may depend on these fixed data and on
fixed cutoffs and compact energy sets. They are independent of the small
parameter $h$ and of the spectral point in the region stated. No uniformity
as $n\downarrow1$ is asserted.

\subsection{The realization and the two spectral parameters}
The realization of the transmission operator used by Petkov and Vodev
\cite[Section~3]{PV} acts on
$L^2(\Omega,n\,dx)\oplus L^2(\Omega,dx)$ with the same maximal graph
domain as $A$. The identity map to $L^2(\Omega)^2$ is a bounded
isomorphism, since $n$ is constant. It conjugates the operators and their
resolvents, so it preserves Riesz projection ranks. We shall use their
counting identity with precisely this multiplicity.

We write $\lambda$ for an energy eigenvalue and use
\begin{equation}\label{eq:scales}
 \zeta=h^2\lambda,\qquad k=\sqrt\lambda,\qquad \Re k\ge0.
\end{equation}
The wave number $k$ is used only to locate the spectrum. The counting
function always counts energies, once each.

\begin{lemma}\label{lem:regularity}
Every eigenfunction of $A$ at a nonzero eigenvalue is smooth up to the
boundary. There is a punctured neighborhood of zero containing no nonzero
eigenvalues. In particular, the nonzero counting function is finite on
bounded intervals of radii.
\end{lemma}
\begin{proof}
Let $(u,v)$ be an eigenfunction at $\lambda\ne0$, and set $w=u-v$.
Then $w,\Delta w\in L^2(\Omega)$ and both generalized Cauchy traces of
$w$ vanish. Zero-Dirichlet elliptic regularity gives $w\in H^2(\Omega)$;
the two zero traces then give $w\in H_0^2(\Omega)$. To apply the regularity
statement at this initial graph regularity, subtract the zero-Dirichlet
solution with the same $(1-\Delta)w$. The remaining homogeneous solution
has zero generalized Dirichlet trace and vanishes by Green's identity
against smooth Dirichlet inverse test functions. The usual regularity
theorem therefore applies. We use the trace and elliptic regularity
statements in \cite[Chapter~11, especially Theorem~11.14]{Grubb}.

The equations imply
\begin{equation}\label{eq:w-reconstruction}
 u=-\frac{\Delta w+\lambda w}{(n-1)\lambda},\qquad
 v=-\frac{\Delta w+n\lambda w}{(n-1)\lambda},
\end{equation}
and hence
\begin{equation}\label{eq:clamped-pencil}
 \Delta^2w+(n+1)\lambda\Delta w+n\lambda^2w=0,
 \qquad w\in H_0^2(\Omega).
\end{equation}
Here $w\ne0$, since $u=v$ would force $(n-1)\lambda u=0$.
The right-hand side of the fourth-order equation is initially in $L^2$.
Clamped elliptic regularity gives $w\in H^4$ and then, by iteration,
$w\in C^\infty(\overline\Omega)$. Formula \eqref{eq:w-reconstruction}
gives the same conclusion for $u,v$.

On $H_0^2(\Omega)$, the norm $\|\Delta w\|_2$ is equivalent to the
$H^2$ norm. Indeed, integration by parts for compactly supported smooth
approximants gives $\sum_{i,j}\|\partial_i\partial_jw\|_2^2
=\|\Delta w\|_2^2$, and Poincare's inequality controls the lower
derivatives. Testing \eqref{eq:clamped-pencil} against $w$ therefore shows
that it has no nonzero solution for sufficiently small $|\lambda|$.
Equivalently, its weak form is an invertible form at $\lambda=0$ plus
$\lambda C_1+\lambda^2C_2$, with $C_1,C_2$ bounded in that Hilbert norm.
Away from zero we use the qualitative discreteness and finite-rank
meromorphic resolvent construction of \cite[Theorem~1.1 and Section~3]{PV}.
This proves the bounded-radius assertion. No compact-resolvent assertion
for $A$ at zero is needed.
\end{proof}

Vodev's eigenvalue-free region applies to the smooth eigenfunctions just
obtained. In the isotropic case, \cite[Theorem~2.1 and the following
remark]{Vdn} gives a constant $C_k$ such that all sufficiently large
transmission wave numbers satisfy
\begin{equation}\label{eq:wave-strip}
 |\Im k|\le C_k.
\end{equation}
The remark is used for $0\le\Re k\le1$; Theorem~2.1 supplies the region
$\Re k>1$. Thus \eqref{eq:wave-strip} covers all large complex energies,
including the exclusion of a large negative-energy branch.

For $m=1,n$, let $G_{m,D}=-m^{-1}\Delta_D$ be the positive Dirichlet
operator, and let $D_m(t)$ count its eigenvalues at most $t$.
The sharp scalar Dirichlet Weyl estimate \cite[(0.1)]{PTL} gives
\begin{equation}\label{eq:scalar-weyl}
 D_m(t)=\tau_m t^{d/2}+O(t^{q/2}),\qquad
 \tau_m=\frac{\omega_d}{(2\pi)^d}|\Omega|m^{d/2}.
\end{equation}
The scaling follows from $D_m(t)=D_1(mt)$. Consequently, in a fixed
compact positive scaled-energy interval, every interval of length $Ch$
has combined scaled Dirichlet multiplicity $O(h^{-q})$, for fixed $C$.
To see this, subtract \eqref{eq:scalar-weyl} at the endpoints divided by
$h^2$. The leading difference and both remainders are $O(h^{-q})$.
Enlarging the interval by $h$ proves the same assertion for closed
endpoints and for a single eigenvalue's multiplicity.

\subsection{Dirichlet-to-Neumann maps and the source approximation}
Let $\Gamma=\partial\Omega$, and use the inner unit normal on $\Gamma$.
For $\zeta$ outside the two scaled Dirichlet spectra, define
\begin{equation}\label{eq:DN}
 \mathcal N_m(\zeta,h)f=h\partial_{\nu_{\rm in}}u|_\Gamma,
 \qquad (h^2\Delta+m\zeta)u=0,\quad u|_\Gamma=f.
\end{equation}
The families extend meromorphically, with simple finite-rank poles at
scaled Dirichlet values. Put
\begin{equation}\label{eq:Lambda}
 \Lambda_h=(I-h^2\Delta_\Gamma)^{1/2},\qquad
 \|f\|_{H_h^1(\Gamma)}=\|\Lambda_hf\|_{L^2(\Gamma)}.
\end{equation}
This norm is uniformly equivalent to the local semiclassical Sobolev
norm. The boundary Laplacian is nonpositive; the square root is its
positive spectral one. Write $r_0(x',\xi')$ for the principal symbol of
$-\Delta_\Gamma$, and $\ell=(1+r_0)^{1/2}$.

We use semiclassical quantization with local phase
$e^{-i\langle x'-y',\xi'\rangle/h}$ and factor $(2\pi h)^{-q}$,
as in \cite{Vdn}. The opposite phase is obtained by
$\xi'\mapsto-\xi'$. It gives the same symbol estimates and the same
principal symbols here, since $r_0$ is even in $\xi'$. We use a linear
partition-of-unity quantization with $\Op_h(1)=I$.

The approximation from \cite[Theorem~1.2]{Vdn} will be used in its full
mapping norm. Work first with a semiclassical parameter $h_e$ such that
$\Re((h_e k)^2)=1$. For every fixed sufficiently small $\delta>0$, let
$\chi_{m,\delta}$ be the source cutoff equal to one in a
$\delta^2$ neighborhood of the glancing set $\{r_0=m\}$ and supported
in a slightly larger such neighborhood. With
\[
 \rho_m=\sqrt{r_0-m(h_e k)^2},\qquad \Re\rho_m<0,
\]
the source gives
\begin{equation}\label{eq:Vodev-interface}
 \left\|h_eN(k;m)-\Op_{h_e}
 \big(\rho_m(1-\chi_{m,\delta})+h_e b\big)\right\|_{L^2\to H_{h_e}^1}
 \le C\delta.
\end{equation}
Here $N(k;m)$ is the unscaled inner-normal map for
$(\Delta+mk^2)u=0$. Here $S_0^0$ denotes symbols of order zero with standard symbol
seminorms uniform in the semiclassical parameter. The symbol $b\in S_0^0$
is independent of both $k$ and the index function. For each $0<\epsilon<1$, the estimate holds when
\[
 C_\delta\le|\Im k|\le(\Re k)^{1-\epsilon},\qquad
 \Re k\ge C_{\epsilon,\delta}.
\]
The domain may be any smooth bounded domain. No concavity is part of
this theorem. The index independence of $b$ and the $H^1$ output in
\eqref{eq:Vodev-interface} are both needed below.

We also use the closed-manifold symbol calculus, boundedness and Garding
estimates, in the conventions of \cite[Chapter~4 and Appendix~E]{EZ}.
For fixed cutoffs, their symbol seminorms are uniform in $h$.
The positive operator in \eqref{eq:Lambda} has symbol $\ell+hS^0$.
One way to justify this for the unbounded square root is to start from
the norm-convergent integral
\begin{equation}\label{eq:square-root-integral}
 (I-h^2\Delta_\Gamma)^{-1/2}
 =\frac2\pi\int_0^\infty(I-h^2\Delta_\Gamma+t^2)^{-1}\,dt.
\end{equation}
The elliptic inverse starts with $(1+r_0+t^2)^{-1}$. Its coefficient at
order $h^j$ has joint order $-2-j$ in $(\xi',t)$, by the usual
composition recursion. Integrating in $t$ gives order $-1-j$, with
principal symbol $\ell^{-1}$. An arbitrarily long recursion makes the
remainder integrable in any fixed Sobolev seminorm; the exact positive
resolvent bounds the correction. Multiplication by
$I-h^2\Delta_\Gamma$ gives the stated symbol of $\Lambda_h$.
In particular, $\Lambda_h:H_h^1\to L^2$ has norm one, and
$\Lambda_h:H^{s+1}\to H^s$ is bounded uniformly for $h\le1$ at each
fixed ordinary Sobolev order.

\section{A normalized boundary determinant}\label{sec:det}

The normalization is chosen so that its principal symbol lies in a fixed
sector and tends to one at high tangential frequency. The first property
will control the logarithm; the second will make its trace-class part
have boundary-order size. Define
\begin{equation}\label{eq:B}
 T_h=\mathcal N_n-\mathcal N_1,\qquad
 B_h(\zeta)=\frac{2}{(n-1)\zeta}\Lambda_hT_h(\zeta).
\end{equation}
For each fixed $h$, this is a bounded meromorphic family on boundary
$L^2$. Its high-frequency order is zero, since the first-order terms of
the two Dirichlet-to-Neumann maps cancel. Fix a compact rectangle
$\mathcal D$ inside
$\{1/2<\Re\zeta<3,\ |\Im\zeta|<9/10\}$, whose interior contains
$[3/5,14/5]$. Every later contour lies in its interior for small $h$.

\subsection{A uniform sector at height \texorpdfstring{$h$}{h}}
\begin{lemma}\label{lem:sector}
There are constants $C_*,c_B,C_B>0$ such that, for every fixed $H>C_*$
and sufficiently small $h$, the following estimates hold uniformly for
$3/5\le x=\Re\zeta\le14/5$ and $C_*h\le|y|=|\Im\zeta|\le Hh$:
\begin{equation}\label{eq:sector-operator}
 \|B_h(\zeta)\|\le C_B,\qquad
 \Re\big(e^{-i\theta}\langle B_h(\zeta)v,v\rangle\big)
 \ge c_B\|v\|_2^2,\quad \theta=\sgn(y)\pi/4.
\end{equation}
The constants $C_*,c_B,C_B$ can be fixed before $H$; the smallness
threshold for $h$ may then depend on $H$.
\end{lemma}
\begin{proof}
Suppose first that $y>0$, and choose
\[
 \alpha_m=\sqrt{r_0-m\zeta},\qquad \Re\alpha_m>0,\quad\Im\alpha_m<0.
\]
The inner-normal principal symbol is $-\alpha_m$. Thus the principal
symbol of \eqref{eq:B} is
\begin{equation}\label{eq:a}
 a=\frac{2\ell(\alpha_1-\alpha_n)}{(n-1)\zeta}
   =\frac{2\ell}{\alpha_n+\alpha_1}.
\end{equation}
The two roots lie in the fourth quadrant, so $a$ lies in the first.
Moreover,
\begin{equation}\label{eq:scalar-sector}
 0<c_0\le|a|\le C_0,\qquad
 \Re(e^{-i\pi/4}a)\ge c_0/\sqrt2.
\end{equation}
For large $r_0$, the denominator is comparable to $2\sqrt{r_0}$.
On a bounded frequency set, the identity
$\alpha_1^2-\alpha_n^2=(n-1)\zeta$ bounds at least one root away
from zero. Vectors in the same quadrant cannot cancel, proving the
remaining bounds in \eqref{eq:scalar-sector}.

For $x$ in the stated compact interval, put $h_e=h/\sqrt x$ in
\eqref{eq:Vodev-interface}. Then $\Re((h_e k)^2)=1$. The change
$\xi'=\sqrt x\eta$ gives
$\Op_{h_e}(c(\eta))=\Op_h(c(\xi'/\sqrt x))$.
Multiplication of \eqref{eq:Vodev-interface} by $\sqrt x$ transforms
its root to $-\alpha_m$ and its common correction to
$hb(x',\xi'/\sqrt x)$. This correction cancels between the two
indices. The two $H^1$ norms are uniformly equivalent. Also
\[
 \Im k=\frac{y}{2h\Re\sqrt\zeta}.
\]
Choosing $C_*$ sufficiently large meets the lower source bound; for
fixed $H$ the upper source bound holds for small $h$, for example with
$\epsilon=1/2$.

On a deleted neighborhood $|r_0-mx|\le C\delta^2$, one has
$|\alpha_m|\le C(\delta+\sqrt{Hh})$. For fixed $\delta$ and small
$h$, deleting it therefore changes the symbol in \eqref{eq:a} by
$O(\delta)$ in supremum norm. The deleted frequencies are bounded.
Outside them the cut symbols have uniform $h$ seminorms, with constants
depending on $\delta$. Their difference has order $-1$ at infinity.
Compose with $\Lambda_h$ and use the actual $L^2\to H_h^1$ error
in \eqref{eq:Vodev-interface}. It becomes an $O(\delta)$ error on
$L^2$. Composition and Garding for the smooth cut symbol give
\[
 \Re(e^{-i\pi/4}\langle B_hv,v\rangle)
 \ge(c_0/\sqrt2-C\delta-C_\delta h)\|v\|_2^2,
\]
as well as the upper norm bound. Fix $\delta$ first and then $h$.
Conjugating the real-coefficient boundary problems gives the lower-half
estimate with angle $-\pi/4$. The auxiliary cut symbols need not be
holomorphic: the estimates concern the exact family $B_h$.
\end{proof}

\subsection{Holomorphic normalization and trace norm}
Let $\chi_h=\Op_h(\chi)$, where $\chi$ equals one for $r_0\le R$
and zero for $r_0\ge2R$. The fixed $R$ is chosen above all glancing
frequencies on $\mathcal D$. The elliptic Poisson construction in
\cite[Section~2, (2.7), (2.9), (2.13)]{PV}, at an arbitrarily large
fixed order $J$, gives a holomorphic approximation $Q_h$ to
$B_h(I-\chi_h)$. We take the actual composition with $\Lambda_h$,
without discarding any further composition remainder, and set
\begin{equation}\label{eq:P}
 P_h=Q_h+\chi_h.
\end{equation}
The local Poisson coordinate in that construction is $x_1>0$ inward.
With $D_{x_1}=-ih\partial_{x_1}$, our inner derivative is
$h\partial_{x_1}=iD_{x_1}$. The source phase has
$\rho=i\sqrt{r_0-m\zeta}$, and hence its actual inner-normal symbol
is $i\rho=-\sqrt{r_0-m\zeta}$. This convention fixes the sign in
\eqref{eq:a}; the common change of normal in the transmission index
identity has no effect.

The leading symbol of $P_h$ is $a(1-\chi)+\chi$. At high frequency
$a=1+O(r_0^{-1})$, with the same decay under the finite derivatives
needed for boundedness. The first subprincipal difference of the two
Poisson boundary symbols has order $-2$ when both leading differential
coefficients are one; see the statement in the proof of
\cite[Lemma~3.3, p.~18, using (2.8)]{PV}. The coefficient at $h^j$, $j\ge2$, has order at
most $1-j$. After $\Lambda_h$ all these terms are bounded and carry
a positive power of $h$. It follows that
\begin{equation}\label{eq:P-close}
 \|P_h-I\|\le\varepsilon_R+C_Rh,\qquad\varepsilon_R\longrightarrow0
 \quad(R\longrightarrow\infty),
\end{equation}
uniformly on $\mathcal D$. Choose $R$ and then $h$ so that this norm
is less than $\varepsilon_0<1/2$, where
$2C_B\varepsilon_0<c_B/2$. The Neumann series makes $P_h^{-1}$
holomorphic on the full rectangle and uniformly bounded by two.

We next keep the smoothing error separate from the low frequencies.
Let
\[
 \delta_D(\zeta,h)=\min\{1,\operatorname{dist}
 (\zeta,h^2\spec G_{1,D}\cup h^2\spec G_{n,D})\}.
\]
The $L^2$ spectral theorem and ordinary zero-Dirichlet regularity give,
for every fixed integer $j\ge1$,
\begin{equation}\label{eq:rough-resolvent}
 \|(h^2G_{m,D}-\zeta)^{-1}f\|_{H^{2j}}
 \le C_jh^{-2j}\delta_D^{-1}\|f\|_{H^{2j-2}}.
\end{equation}
Indeed the $L^2$ norm is at most $\delta_D^{-1}$, and iteration of
\[
 G_{m,D}(h^2G_{m,D}-\zeta)^{-1}f
 =h^{-2}\big(f+\zeta(h^2G_{m,D}-\zeta)^{-1}f\big)
\]
costs $h^{-2}$ at each step of the elliptic estimate. It does not add
another distance factor. In particular, this argument does not require
a higher ordinary-Sobolev resolvent bound uniform in $h$.\footnote{We
do not use (2.3) of \cite{PV}. Its displayed $h$ factor fails for the
single Dirichlet-to-Neumann map on the unit three-ball at scaled energy
$-1$ with boundary datum $1$: the solution is
$\sinh(r/h)/(r\sinh(1/h))$, and the convention
$-ih\partial_{\nu_{\rm out}}$ gives
$-i(\coth(1/h)-h)$. The truncated distance to the positive Dirichlet
spectrum is one, but this value does not tend to zero. Only the pure
Poisson construction, before that auxiliary estimate, is used here.}

The exact corrected Poisson solution gives
\begin{equation}\label{eq:exact-residual}
 B_h-P_h=(B_h-I)\chi_h+E_{J,h}.
\end{equation}
Fix an integer $s>q/2$. The pure interior error in
\cite[(2.13)]{PV}, with its index $s+1$, maps boundary $L^2$ to
$H^{2s+7/2}(\Omega)$ with norm
$O(h^{J-2s-2-\ell_1})$, where $\ell_1$ is independent of $s,J$ and
$s+1\le J/4$. Embed into $H^{2s+2}$ and use
\eqref{eq:rough-resolvent} with $j=s+2$. The result is in
$H^{2s+4}(\Omega)$, at a cost $h^{-2s-4}\delta_D^{-1}$.
The scaled normal trace has output $H^{2s+5/2}(\Gamma)$ and adds a
factor $h$. Finally $\Lambda_h$ maps to $H^{2s+3/2}(\Gamma)$, which
embeds into $H^{2s}(\Gamma)$. Thus
\begin{equation}\label{eq:smoothing-error}
 \|E_{J,h}\|_{L^2\to H^{2s}}
 \le C_{s,J}h^{J-4s-5-\ell_1}\delta_D^{-1}.
\end{equation}
Choose $J\ge4(s+1)$ and $J>4s+6+\ell_1$. At height $|y|\ge C_*h$,
the right-hand side is bounded. The embedding $H^{2s}(\Gamma)\to L^2$
is trace class, since $2s>q$.

The low-frequency cutoff has
\begin{equation}\label{eq:cutoff-trace}
 \|\chi_h\|_1=O(h^{-q}).
\end{equation}
For completeness, its norm on $L^2$ is $O(1)$ and its norm into
$H^{2s}$ is $O(h^{-2s})$. The singular values of the fixed embedding
are $O(j^{-2s/q})$; see \cite[(2.11)]{PV} or
\cite[Appendix~E]{EZ}. Hence
$s_j(\chi_h)\le C_s\min(1,h^{-2s}j^{-2s/q})$.
Summing below and above $j\asymp h^{-q}$ proves
\eqref{eq:cutoff-trace}.

Combining \cref{lem:sector} with \eqref{eq:exact-residual}--
\eqref{eq:cutoff-trace}, define
\begin{equation}\label{eq:Fredholm}
 F_h=B_hP_h^{-1}=I+K_h,\qquad g_h=\det(I+K_h),\qquad M=h^{-q}.
\end{equation}
On the critical strips $C_*h\le|y|\le Hh$,
\begin{equation}\label{eq:trace-sector}
 \|K_h\|_1\le CM,\qquad
 \Re\big(e^{-i\theta}\langle F_hv,v\rangle\big)
 \ge(c_B/2)\|v\|_2^2.
\end{equation}
For the second estimate use
$\|F_h-B_h\|\le2C_B\varepsilon_0<c_B/2$.
The family is meromorphic in trace norm throughout $\mathcal D$,
with finite-rank principal parts. We shall also use a coarse bound
\begin{equation}\label{eq:coarse}
 \log|g_h(\zeta)|\le\|K_h(\zeta)\|_1
 \le Ch^{-A_0}\delta_D^{-1}
\end{equation}
for one fixed finite $A_0$. To obtain it, apply a fixed smooth boundary
extension to $\chi_hf$, subtract the Dirichlet resolvent of its equation
residual, and use \eqref{eq:rough-resolvent} at a fixed sufficiently
large order. Cutoff differentiation and each fixed trace cost only
finitely many powers of $h^{-1}$. The high-frequency remainder is
already covered by \eqref{eq:smoothing-error}. Factoring through the
trace-class Sobolev embedding gives \eqref{eq:coarse}. Only this
polynomial form will be needed at remote vertical edges.

\subsection{A complex logarithm and the divisor identity}
The trace norm in \eqref{eq:trace-sector} controls a particular analytic
logarithm, including its imaginary part.
\begin{lemma}\label{lem:log}
On each of the upper and lower critical strips there is a holomorphic
function $L_h^\pm$ such that
\begin{equation}\label{eq:canonical-log}
 e^{L_h^\pm}=g_h,\qquad |L_h^\pm|\le CM,\qquad
 (L_h^\pm)'=g_h'/g_h.
\end{equation}
\end{lemma}
\begin{proof}
For $0\le t\le1$, $I+tK_h=(1-t)I+tF_h$. Its rotated real part is
bounded below by $\min(\cos\theta,c_B/2)\|v\|_2^2$.
This gives injectivity and closed range; the corresponding conjugate
rotation for the adjoint gives dense range. Thus its inverse is
uniformly bounded. Set
\[
 L_h^\pm(\zeta)=\int_0^1\Tr\big((I+tK_h)^{-1}K_h\big)\,dt.
\]
The integrand is holomorphic and uniformly bounded in trace norm by
$CM$. The determinant differentiation formula, integrated from $t=0$
where the determinant is one, proves \eqref{eq:canonical-log}.
These trace-class formulas can also be obtained directly by finite-rank
approximation: the exterior-power determinant series is absolutely
bounded by $e^{\|K_h\|_1}$, converges locally uniformly in trace norm,
and has the finite-dimensional differentiation formula where invertible.
Cauchy's formula passes its derivatives to the limit.
\end{proof}

For a scaled Dirichlet value $p$, let $m_D(p)$ be the sum of its two
Dirichlet multiplicities. Let $m_T(p)$ be the Riesz multiplicity of
$p/h^2$ for $A$, or zero if it is not a transmission eigenvalue.
Away from the Dirichlet spectra put $m_D=0$ as well. Orders of zeros
are positive and orders of poles negative.

\begin{proposition}\label{prop:index}
The determinant in \eqref{eq:Fredholm} satisfies
\begin{equation}\label{eq:divisor}
 \ord_p g_h=m_T(p)-m_D(p).
\end{equation}
Equivalently, for every simple positively oriented contour $\gamma$
in $\mathcal D$ avoiding the transmission and Dirichlet spectra,
\begin{equation}\label{eq:index}
 \#_T(\gamma)=\#_{D,1}(\gamma)+\#_{D,n}(\gamma)
 +\frac1{2\pi i}\int_\gamma\frac{g_h'}{g_h}\,d\zeta.
\end{equation}
All transmission counts here use the original Riesz ranks, including at
common Dirichlet values.
\end{proposition}
\begin{proof}
The original-operator trace identity is
\cite[Proposition~3.1, Corollary~3.2 and Lemma~3.6]{PV}. Its
meromorphic Fredholm hypothesis holds here by
\eqref{eq:exact-residual} and the invertible critical-strip values in
\eqref{eq:trace-sector}. The identity between the constant-weighted source Hilbert space and our
unweighted $L^2$ space is a bounded conjugacy, so it preserves the
resolvents and Riesz ranks; see \cref{sec:prelim}.

It remains to check the normalizations. Passing to $B_h$ multiplies
the boundary difference on the left by a holomorphic nonzero scalar,
a fixed common normal sign and the invertible map
$\Lambda_h:H_h^1\to L^2$. Passing to $F_h$ multiplies it on the right
by the holomorphic invertible $P_h^{-1}$. None changes the local
meromorphic index. Indeed, finite-dimensional projections separate an
invertible complementary block from the finite singular block of a
Fredholm family. Holomorphic invertible changes multiply its Schur
complement by holomorphic invertible factors, whose determinant orders
are zero. For an identity-plus-trace-class representative the remaining
order is that of its Fredholm determinant. The source identity therefore
gives \eqref{eq:index}, and small circles give \eqref{eq:divisor}.
In particular, pole order is at most $m_D(p)$, even when a transmission
point coincides with one or both Dirichlet spectra.
\end{proof}

\section{Pole cancellation in a thin strip}\label{sec:strip}

We now use $g_h=\det(I+K_h)$ from \eqref{eq:Fredholm} and
$M=h^{1-d}$. Its only possible poles are real Dirichlet values; its
upper and lower logarithms are bounded by $CM$ at imaginary height
comparable to $h$. The purpose of the pole factors is to use those two
bounds across the real axis.

Fix
\begin{equation}\label{eq:heights}
 L>\max\{2C_*,8(C_k+1)\},\qquad H>20L.
\end{equation}
These constants are chosen before reducing $h$. Recall from
\eqref{eq:wave-strip} that $|\Im\sqrt\lambda|\le C_k$ and
$\zeta=h^2\lambda$. In the band $\Re\zeta\le3$ this gives
\[
 |\Im\zeta|\le2C_k\sqrt{3+C_k^2h^2}\,h<(L/2)h
\]
for all sufficiently large eigenvalues and small $h$.
Choose $a_h\in[3/5,7/10]$ and $b_h\in[27/10,14/5]$ at distance
at least $h^{d+2}$ from the scaled Dirichlet spectra. Such choices exist:
there are $O(h^{-d})$ relevant Dirichlet values counted with multiplicity,
so the excluded intervals have total length $O(h^2)$. Set
\[
 S_h=\{a_h<\Re\zeta<b_h,\ |\Im\zeta|<Hh\},\qquad
 b_p(\zeta)=\tanh\frac{\pi(\zeta-p)}{4Hh},
\]
and form the finite product
\begin{equation}\label{eq:pole-product}
 \mathcal B_h(\zeta)=\prod_{a_h<p<b_h}b_p(\zeta)^{m_D(p)},\qquad
 f_h=g_h\mathcal B_h.
\end{equation}
Each factor is analytic in the strip, has just the simple zero $p$
there, and satisfies $|b_p|\le1$ inside and $|b_p|=1$ on its horizontal
edges. These assertions follow from
\begin{equation}\label{eq:tanh-modulus}
 |\tanh(u+iv)|^2=
 \frac{\sinh^2u+\sin^2v}{\sinh^2u+\cos^2v},\qquad |v|\le\pi/4.
\end{equation}
The nearest poles are at $p\pm2iHh$. By \eqref{eq:divisor}, $f_h$
is holomorphic on a neighborhood of the closed rectangle. Its zero
orders at all transmission points, including common Dirichlet points,
are exactly $m_T$.

\subsection{Upper bounds and local counts}
On the horizontal edges \cref{lem:log} gives
$\log|f_h|\le CM$. On the vertical edges \eqref{eq:coarse} and the
chosen distance give the coarser bound $Ch^{-A_0-d-2}$.
This large bound has exponentially small influence in the interior.
For a direct comparison, put $\kappa=\pi/(4Hh)$ and
\[
 v_h(x+iy)=\sqrt2\big(e^{-\kappa(x-a_h)}
                  +e^{-\kappa(b_h-x)}\big)\cos(\kappa y).
\]
It is harmonic and nonnegative on $S_h$, and is at least one on both
vertical edges. Therefore $CM+Ch^{-A_0-d-2}v_h$ majorizes
$\log|f_h|$. To justify this even at zeros, take a holomorphic function
whose real part is the harmonic majorant and apply the maximum-modulus
principle to $f_h$ times its negative exponential. If
$4/5\le x\le11/5$, then $v_h\le Ce^{-c/h}$. Hence
\begin{equation}\label{eq:interior-upper}
 \log|f_h(\zeta)|\le CM,
 \qquad 4/5\le\Re\zeta\le11/5,\quad |\Im\zeta|\le Hh.
\end{equation}

For $a$ in a fixed compact inner interval containing $[1,2]$, let
$z_a=a+iLh$. The lower bound in \eqref{eq:canonical-log} gives
$\log|g_h(z_a)|\ge-CM$. Formula \eqref{eq:tanh-modulus} also gives
\begin{equation}\label{eq:pole-weight}
 0\le-\log|b_p(z_a)|\le Ce^{-c|a-p|/h}.
\end{equation}
For bounded $|a-p|/h$, the fixed positive imaginary argument
$\pi L/(4H)$ bounds the modulus away from zero. For larger separation,
use the exponential convergence of $\tanh$ to $1$ or $-1$.
Each interval of length $h$ has combined scaled Dirichlet multiplicity
$O(M)$, by \eqref{eq:scalar-weyl}. Partitioning the real axis into
these intervals and summing the exponential weights yields
\begin{equation}\label{eq:center-lower}
 -\log|\mathcal B_h(z_a)|\le CM,\qquad
 \log|f_h(z_a)|\ge-CM.
\end{equation}
The constants are uniform in $a$.

For small $h$, the disk of radius $8Lh$ about $z_a$ lies in the region
of \eqref{eq:interior-upper}. Jensen's formula and
\eqref{eq:center-lower} bound the number of its zeros in radius $7Lh$
by $CM$: each such zero contributes at least $\log(8/7)$ to the
Jensen sum. If the outer circle contains zeros, use nearby radii tending
to $8Lh$. The smaller disk $|\zeta-a|\le2Lh$ lies inside that disk.
Since the zeros of $f_h$ have precisely the transmission orders, we obtain
\begin{equation}\label{eq:local-TE}
 \sum_{|\zeta-a|\le2Lh}m_T(\zeta)\le CM.
\end{equation}
In particular, this bound holds at every real endpoint, without a
spacing assumption. The elementary disk formulas used here and below
follow from Cauchy's formula, the maximum principle and Schwarz's lemma;
see \cite[Sections~6--8, 13 and 17]{Bonk}.

\subsection{The argument on a short vertical segment}
\begin{lemma}\label{lem:vertical}
For every such $a$, there is an $a'$ with $|a'-a|<h/10$ such that
the segment $\gamma_{a'}=[a'-iLh,a'+iLh]$ avoids both spectra and
\begin{equation}\label{eq:vertical-argument}
 \left|\Im\int_{\gamma_{a'}}\frac{g_h'}{g_h}\,d\zeta\right|\le CM.
\end{equation}
\end{lemma}
\begin{proof}
Choose $R_a\in[6Lh,7Lh]$ so that its circle about $z_a$ contains no
zero of $f_h$. Factor the zeros in this disk by its disk Blaschke
product $\mathcal Z_h$, and write
\[
 f_h=\mathcal Z_h e^{w_h}.
\]
The quotient is zero-free on a neighborhood of the closed disk, so its
holomorphic logarithm $w_h$ exists there. There are $O(M)$ factors.
They have unit modulus on the circle and modulus at most one inside.
Thus \eqref{eq:interior-upper} and \eqref{eq:center-lower} imply
\[
 \Re w_h\le CM,\qquad \Re w_h(z_a)\ge-CM.
\]
The usual real-part form of the Cauchy estimate gives
\begin{equation}\label{eq:zero-free-derivative}
 |w_h'(\zeta)|\le CM/h\quad\text{for }|\zeta-z_a|\le4Lh.
\end{equation}
For clarity, it follows by setting
$W=U-w_h+i\Im w_h(z_a)$ with $U\ge\sup\Re w_h$.
If $W(z_a)>0$, the function
$(W-W(z_a))/(W+W(z_a))$ maps the disk to the unit disk and vanishes
at its center. Schwarz's lemma bounds $W$ on a smaller disk by a
constant times $W(z_a)$, and Cauchy's estimate bounds its derivative.
If $W(z_a)=0$, the real-part minimum principle makes it constant.
The fixed radius ratios give \eqref{eq:zero-free-derivative}.

Choose $|a'-a|<h/10$ avoiding the finitely many real parts of local
zeros and all real Dirichlet values. The resulting segment is inside
$|\zeta-z_a|<4Lh$. For a linear zero factor $\zeta-w$, the change of argument is the
angle subtended by a straight segment, so
\[
 \left|\Im\int_{\gamma_{a'}}\frac{d\zeta}{\zeta-w}\right|\le\pi.
\] In the coordinate
$z=\zeta-z_a$, a disk factor has denominator
$R_a^2-\overline{w-z_a}\,z$; on the inner disk its logarithmic
derivative is $O(h^{-1})$, uniformly in the zero. Its integral is
therefore $O(1)$. Summing the $O(M)$ factors and
\eqref{eq:zero-free-derivative} proves
\[
 \left|\Im\int_{\gamma_{a'}}f_h'/f_h\,d\zeta\right|\le CM.
\]

It remains to subtract $\mathcal B_h'/\mathcal B_h$.
For $|p-a|\le10Hh$, the difference
$b_p'/b_p-(\zeta-p)^{-1}$ is $O(h^{-1})$ on the segment,
uniformly even as $a'-p$ tends to zero. This follows by scaling by
$h$ and removing its simple pole; the other zeros and poles lie
outside the strip. There are $O(M)$ near factors. Their removed
linear terms contribute at most $\pi$ each. For the remaining factors,
\[
 |b_p'/b_p|\le Ch^{-1}e^{-c|a-p|/h}.
\]
The length-$h$ Dirichlet bound and summation of this exponential tail,
followed by integration over length $2Lh$, give $O(M)$ again.
Using \eqref{eq:pole-product} proves \eqref{eq:vertical-argument}.
\end{proof}

\section{The Weyl formula}\label{sec:count}

We can now compare the transmission and scalar Dirichlet counts with
boundary-order error. Apply \cref{lem:vertical} at $1$ and $2$ to
choose $a'=1+O(h)$ and $b'=2+O(h)$, with shifts less than $h/10$.
Let $\mathcal R_h$ be the positively oriented rectangle with these
vertical sides and horizontal sides at $\pm Lh$. Its boundary avoids
the two spectra, including common Dirichlet/transmission points.
On a horizontal side,
\[
 \int g_h'/g_h\,d\zeta=L_h^\pm(\text{end})-L_h^\pm(\text{start})
 =O(M)
\]
by the actual analytic branch in \cref{lem:log}. The two short vertical
sides satisfy \eqref{eq:vertical-argument}. Since the complete contour
integral divided by $2\pi i$ is real, these imaginary-part estimates
and \cref{prop:index} give
\begin{equation}\label{eq:rect-count}
 \#_T(\mathcal R_h)-\#_{D,1}(\mathcal R_h)-\#_{D,n}(\mathcal R_h)
 =O(h^{1-d}).
\end{equation}

For every sufficiently large transmission energy in a fixed positive
scaled band, \eqref{eq:wave-strip} implies
\begin{equation}\label{eq:radial-real}
 |\Im\zeta|=O(h),\qquad
 |\zeta|-\Re\zeta=
 \frac{(\Im\zeta)^2}{|\zeta|+\Re\zeta}=O(h^2).
\end{equation}
The choice \eqref{eq:heights} places all of them below $Lh/2$.
Replacing $\mathcal R_h$ by the exact annulus
$1<|\zeta|\le2$ can therefore change only points inside the
$2Lh$ disks about $1$ and $2$. Their total transmission multiplicity
is $O(M)$ by \eqref{eq:local-TE}. This includes values lying exactly
on a counting circle. The scalar endpoint contribution is also $O(M)$ by
\eqref{eq:scalar-weyl}. Thus \eqref{eq:rect-count} becomes
\begin{equation}\label{eq:annular}
 N(\sqrt2/h)-N(1/h)
 =(2^{d/2}-1)(\tau_1+\tau_n)h^{-d}+O(h^{1-d}).
\end{equation}
Here $\tau_1+\tau_n=\omega_d(2\pi)^{-d}|\Omega|(1+n^{d/2})$ is
the combined scalar Weyl coefficient. Every sufficiently small real $h$ is allowed: the auxiliary vertical
sides were chosen for each $h$, and the local count removes their
effect at the exact radii.

\begin{proof}[Proof of \cref{thm:main}]
Put $c=\tau_1+\tau_n$ and $h=\sqrt2/r$ in \eqref{eq:annular}.
For every sufficiently large real $r$,
\[
 N(r)-N(r/\sqrt2)=(1-2^{-d/2})cr^d+O(r^{d-1}).
\]
Sum over $r,r/\sqrt2,r/2,\ldots$ until the last radius belongs to
a fixed bounded interval. The leading terms telescope, and the error
sum is bounded by
\[
 Cr^{d-1}\sum_{j\ge0}2^{-j(d-1)/2}=O(r^{d-1}),
\]
since $d\ge2$. The remaining bounded-radius count is finite by
\cref{lem:regularity}. Its contribution, and its leading volume term,
are $O(1)$. Substitution of the two constants in
\eqref{eq:scalar-weyl} proves the stated formula.
\end{proof}

\appendix
\section{Common Dirichlet values in a ball}\label{app:ball}

The original algebraic multiplicity matters even in a simple radial
example. Consider the unit ball in $\mathbb R^3$ with $n=4$, restricted
to radial fields. The map $(u,v)\mapsto(f,g)=(ru,rv)$ is unitary up to
a fixed scalar, and the radial operator is
\[
 \mathcal A(f,g)=(-\tfrac14 f'',-g''),\qquad
 f(0)=g(0)=0,\quad f(1)=g(1),\quad f'(1)=g'(1),
\]
on $H^2(0,1)^2$. The zero condition at the origin excludes the
$1/r$ term, whose Laplacian has a delta mass and is not in the original
graph domain. The radial subspace is a reducing block, since rotational
averaging commutes with the operator and traces. This finite-interval
realization has compact resolvent.

At energy $\lambda=k^2\ne0$, its solutions have the form
$f=a\sin(2kr)$ and $g=b\sin(kr)$. The boundary matrix is
\begin{equation}\label{eq:ball-matrix}
 \begin{pmatrix}\sin2k&-\sin k\\2k\cos2k&-k\cos k\end{pmatrix},
 \qquad\det=-2k\sin^3k.
\end{equation}
Compare the corresponding model determinant in
\cite[Section~2]{PS}. Thus the nonzero energies are
$\lambda_m=(m\pi)^2$, $m\ge1$, and each eigenspace has dimension one.
We check its Riesz multiplicity directly rather than identify it merely
with a determinant root order.

Near $\lambda_0=k_0^2$, $k_0=m\pi>0$, choose the analytic root $k(\lambda)$.
The analytic pair
\[
 W(\lambda,r)=\left(\sin(2k(\lambda)r),
 \frac{2\cos(2k(\lambda))}{\cos k(\lambda)}\sin(k(\lambda)r)\right)
\]
solves the differential equations and all but the value-matching
condition. Its value mismatch is
$2\sin^3 k(\lambda)/\cos k(\lambda)$, of exact order three.
The family is an identity for the differential expression, not an
assertion that $W(\lambda)$ lies in the transmission domain for all
$\lambda$. Its Taylor coefficients $W_0,W_1,W_2$ do lie in that
domain and satisfy
\[
 (\mathcal A-\lambda_0)W_0=0,\quad
 (\mathcal A-\lambda_0)W_1=W_0,\quad
 (\mathcal A-\lambda_0)W_2=W_1.
\]
There is no chain of length $\ell\ge4$. Its formal Taylor polynomial
would solve the two initial-value ODEs modulo
$(\lambda-\lambda_0)^\ell$, hence have the form
$a(\lambda)\sin(2k(\lambda)r),b(\lambda)\sin(k(\lambda)r)$ modulo
that order. The leading eigenvector has $a(\lambda_0)\ne0$.
Derivative matching forces
$b=2a\cos2k/\cos k$; value matching would then require
$2a\sin^3k/\cos k=O((\lambda-\lambda_0)^\ell)$, a contradiction.
The one-dimensional eigenspace therefore belongs to one Jordan block
of size three. Its original Riesz multiplicity is three.

For constant boundary values, the outward DN difference is
\[
 2k\cot2k-k\cot k=-k\tan k.
\]
At $k=m\pi$ its order is one, while $m_D=2$ and $m_T=3$;
thus $1=3-2$ in \eqref{eq:divisor}. At a nonzero half-integer
multiple of $\pi$ it has a simple pole and $-1=0-1$.
The exact radial count is
\[
 N_{\rm rad}(r)=3\lfloor r/\pi\rfloor,
\]
and the sum of the radial Dirichlet counts is
$\lfloor r/\pi\rfloor+\lfloor2r/\pi\rfloor$. Their difference is
always zero or minus one, including thresholds. This calculation concerns
only the radial block; it is not the full three-dimensional volume count
and is not used in the proof of \cref{thm:main}.

\begin{thebibliography}{8}

\bibitem{Bonk}
M.~Bonk, \emph{Complex Analysis}, course notes for Math~246A and 246B,
University of California, Los Angeles,
\url{https://www.math.ucla.edu/~mbonk/complana.pdf}.

\bibitem{EZ}
L.~C. Evans and M.~Zworski,
\emph{Lectures on Semiclassical Analysis}, version~0.2,
University of California, Berkeley,
\url{https://math.mit.edu/~vwg/semiclassicalEvansZworski.pdf}.

\bibitem{Grubb}
G.~Grubb, \emph{Distributions and Operators},
Chapter~11, Pseudodifferential methods for boundary value problems,
author-hosted chapter excerpt,
\url{https://web.math.ku.dk/~grubb/dist11n.pdf}.

\bibitem{PV}
V.~Petkov and G.~Vodev,
Asymptotics of the number of the interior transmission eigenvalues,
\emph{J. Spectr. Theory} \textbf{7} (2017), 1--31,
\href{https://doi.org/10.4171/JST/154}{doi:10.4171/JST/154}.

\bibitem{PS}
H.~Pham and P.~Stefanov,
\emph{Weyl asymptotics of the transmission eigenvalues for a constant
index of refraction}, arXiv:1309.3616v2,
\url{https://arxiv.org/abs/1309.3616v2}.

\bibitem{PTL}
Pham The Lai,
Meilleures estimations asymptotiques des restes de la fonction spectrale
et des valeurs propres relatifs au laplacien,
\emph{Math. Scand.} \textbf{48} (1981), 5--38,
\href{https://doi.org/10.7146/math.scand.a-11895}{doi:10.7146/math.scand.a-11895}.

\bibitem{Vsurvey}
G.~Vodev, Asymptotic behavior of the transmission eigenvalues,
\emph{RIMS K\^oky\^uroku} \textbf{2045} (2017), 61--66,
\url{https://www.kurims.kyoto-u.ac.jp/~kyodo/kokyuroku/contents/pdf/2045-04.pdf}.

\bibitem{Vdn}
G.~Vodev,
High-frequency approximation of the interior Dirichlet-to-Neumann map
and applications to the transmission eigenvalues,
\emph{Anal. PDE} \textbf{11} (2018), 213--236,
\href{https://doi.org/10.2140/apde.2018.11.213}{doi:10.2140/apde.2018.11.213}.

\end{thebibliography}

\section*{AI-use disclosure}
OpenAI Codex was used extensively in the literature and source review,
development and checking of the arguments, and preparation of this
manuscript. One reader simulation and one referee-style review were
conducted in fresh AI contexts on a shared host, and their feedback informed
one revision. These uses do not constitute independent human verification
or formal proof certification. Independent mathematical review is requested.

\end{document}
